% BHU PYQ -- Transcribed & Corrected
% Paper: Paper Code & Name: STAMN-11 Descriptive Statistics
% Exam: B.Sc. Sem- I Mid Semester Examination 2024-25 (NEP Minor)
% Department: Department of Statistics

\documentclass[11pt,a4paper]{article}
\usepackage[a4paper,top=2.5cm,bottom=2.5cm,left=2.8cm,right=2.8cm,headheight=14pt]{geometry}
\usepackage{amsmath,amssymb}
\usepackage[T1]{fontenc}
\usepackage[utf8]{inputenc}
\usepackage{lmodern,microtype,fancyhdr,enumitem,hyperref,lastpage,parskip}

\hypersetup{colorlinks=true,urlcolor=black,linkcolor=black,pdftitle={STAMN11: Descriptive Statistics Mid Term -- BHU Sem I 2024-25 (NEP Minor)}}
\pagestyle{fancy}\fancyhf{}
\fancyhead[L]{\small   \textit{Banaras Hindu University}}
\fancyhead[R]{\small   \textit{STAMN-11: Descriptive Statistics (Mid Term)}}
\fancyfoot[C]{\small Page  \thepage\ of \pageref{LastPage}}


\newcommand{\pts}[1]{\hfill   \textit{[#1]}}


\begin{document}


\begin{center}
  {\large\bfseries B.Sc. Sem- I Mid Semester Examination 2024-25}\\[4pt]
  {\large\bfseries Paper Code \& Name: STAMN-11 Descriptive Statistics}\\[4pt]
  {\normalsize Time: One Hour \hfill Full Marks: 20}
\end{center}

\rule{\linewidth}{0.4pt}\smallskip



\noindent   \textit{Note: Attempt any two questions. Each question carries equal marks.}
\bigskip


\begin{enumerate[label=   \textbf{\arabic*.}, leftmargin=*]
  \item Briefly explain the different types of measurement scales with suitable example.\pts{10}

  \medskip\rule{\linewidth}{0.2pt}\medskip

  \item Find the mean and median of the squares of the first 50 natural numbers.\pts{10}

  \medskip\rule{\linewidth}{0.2pt}\medskip

  \item What do you understand by skewness and kurtosis? Discuss in detail.\pts{10}
\end{enumerate}

\vfill
\rule{\linewidth}{0.4pt}

\begin{center}\small
\href{https://bhu-exam.vercel.app/}{Click here to visit our website}\\[4pt]
\href{https://me-aryan.vercel.app/}{Click here to visit our contributor}
\end{center}

\end{document}
